% Hexagon note, 23 Nov 2026. CC0 1.0.
% Mathematics: GPT-6 Astra (ChatGPT Pro mode). Text: Claude Opus 5.5 at effort high, from the run's text.
% Checks: Claude Opus 5.5 at effort high (targeted proof check, every invariant recomputed).
\documentclass[11pt,a4paper]{article}
\usepackage[T1]{fontenc}
\usepackage{lmodern}
\usepackage{microtype}
\usepackage[margin=2.7cm]{geometry}
\usepackage{amsmath,amssymb,amsthm}
\PassOptionsToPackage{hyphens}{url}
\usepackage[hidelinks]{hyperref}
\emergencystretch=2em

\theoremstyle{plain}
\newtheorem{theorem}{Theorem}
\newtheorem{lemma}[theorem]{Lemma}
\newtheorem{proposition}[theorem]{Proposition}
\newtheorem{corollary}[theorem]{Corollary}
\theoremstyle{remark}
\newtheorem{remark}[theorem]{Remark}

\newcommand{\OO}{\mathcal O}
\newcommand{\F}{\mathbf F}
\newcommand{\PP}{\mathbf P}
\newcommand{\Sym}{\operatorname{Sym}}
\newcommand{\Spec}{\operatorname{Spec}}
\newcommand{\Pic}{\operatorname{Pic}}
\newcommand{\workbenchurl}{https://github.com/KokunoYumeto/mathematics-workbenches/blob/main/works/WB-0003.md}

\title{A very ample counterexample to Kodaira vanishing on a surface,\\ answering a question of Raynaud}
\author{}
\date{23 November 2026 \quad CC0 1.0}

\begin{document}
\maketitle

\begin{abstract}
\noindent
In 1978 Raynaud constructed a smooth projective surface in positive characteristic with an ample line bundle whose
inverse has nonzero first cohomology, and asked whether Kodaira vanishing can also fail for an ample line bundle that
is generated by its global sections, or even very ample. We answer both questions on surfaces, the smallest dimension
in which they have content. On Raynaud's own surface in characteristic~$2$ a modified twist $Q$ is ample and globally
generated with $h^1(Q^{-1})=1$. Taking ninth roots of a general basis of the sections of $Q^{\otimes 9}$ then gives,
over every algebraically closed field of characteristic~$2$, a smooth projective surface $Y$ with a very ample line
bundle $L$ such that $h^0(L)=562$, $h^1(L^{-1})=1$ and $h^1(\OO_Y)=565$. The passage from globally generated to very
ample is general: in any characteristic $p$, roots of degree prime to $p$ of a general basis turn a globally generated
counterexample into a very ample one of the same dimension. The earlier very ample counterexamples, on homogeneous
spaces (Lauritzen) and on varieties of dimension at least six (Lauritzen and Rao), are not surfaces. The mathematics was produced by GPT-6 Astra (ChatGPT Pro mode) and
checked by Claude Opus 5.5 at effort high, which also wrote this text. No human has checked it.
\end{abstract}

\section{Raynaud's question and the result}

Let $X$ be a smooth projective variety of dimension $n$ over an algebraically closed field $k$, and $L$ an ample line
bundle on $X$. Kodaira's vanishing theorem says that $H^i(X,L^{-1})=0$ for $i<n$ when $k$ has characteristic zero. On
a curve the only assertion is $H^0(X,L^{-1})=0$, which holds in every characteristic because $L^{-1}$ has negative
degree. Surfaces are therefore the first case in which the theorem has content, and there the assertion is
$H^1(X,L^{-1})=0$.

Raynaud showed that this fails in every positive characteristic \cite[p.~273]{Raynaud}. His surface $X$ is fibred over a
curve with integral fibres carrying a horizontal line of cusps. In characteristic~$2$ it is a triple cover of a ruled
surface over an Artin--Schreier curve, and the ample bundle is $\mathcal I=\OO_X(E)\otimes N$, where $E$ is a section of
the fibration and $N$ comes from the base. In his \emph{Remarque~3} \cite[p.~278]{Raynaud} he takes, in
characteristic~$2$, the curve $y^2+y=x^{15}$ and $N=\OO_C(2\infty)$. He notes that $\mathcal I$ is generated by its
sections outside $E$, and asks:
\begin{quote}
\guillemotleft\ Peut-on contre-exempler le th\'eor\`eme de Kodaira avec un faisceau ample, engendr\'e par ses
sections, voire tr\`es ample?\ \guillemotright
\end{quote}
Lauritzen answered the very ample question in higher dimension, on homogeneous spaces, by modular representation
theory; Lauritzen and Rao then gave elementary very ample examples of dimension at least six
\cite[pp.~1--2]{LauritzenRao}\cite{Lauritzen}. They point out that their failure is in the penultimate degree and need
not be inherited by hyperplane sections, and that examples of smaller dimension would be welcome
\cite[p.~2]{LauritzenRao}. We answer both of Raynaud's questions for surfaces.

\begin{theorem}\label{thm:main}
Let $k$ be an algebraically closed field of characteristic $2$, for instance $\overline{\F}_2$.
\begin{enumerate}
\item[\textup{(a)}] On Raynaud's surface $S$ over $k$ with the data of his Remarque~3 there is an ample line bundle $Q$,
generated by its global sections, with
\[
 h^0(S,Q)=6,\qquad h^1(S,Q^{-1})=1,\qquad H^1(S,Q^{-n})=0\ \ (n\ge2),\qquad h^1(S,\OO_S)=9 .
\]
\item[\textup{(b)}] There is a smooth projective connected surface $Y$ over $k$, finite and flat of degree $9^{556}$ over
$S$, with a very ample line bundle $L$ such that
\[
 h^0(Y,L)=562,\qquad h^1(Y,L^{-1})=1,\qquad h^1(Y,\OO_Y)=565 .
\]
\end{enumerate}
For every smooth connected curve $H\in|L|$, the kernel of the restriction $\Pic_{Y/k}\to\Pic_{H/k}$ has a
one-dimensional tangent space at the identity.
\end{theorem}

So Kodaira vanishing fails on a surface for a very ample line bundle, already over $\overline{\F}_2$. In particular
$Y$ is defined over a finite field. The last assertion concerns tangent spaces, which is all that $H^1(L^{-1})$
controls. Picard schemes in positive characteristic need not be reduced, and the theorem says nothing about geometric
points of that kernel.

Part (b) follows from part (a) by a construction that works in any dimension and any characteristic.

\begin{theorem}\label{thm:roots}
Let $X$ be a smooth projective connected variety over an algebraically closed field $k$ of characteristic $p\ge0$, and
$Q$ a line bundle on $X$ generated by its global sections such that $Q^{\otimes m}$ is very ample for an integer $m\ge2$
prime to $p$. Let $s_1,\dots,s_\ell$ be a general basis of $H^0(X,Q^{\otimes m})$ and let $\rho:Y\to X$ be the cover
obtained by adjoining an $m$-th root of every $s_i$. Then $Y$ is smooth, projective and connected, $L=\rho^*Q$ is very
ample, and for every $i$
\[
 H^i(Y,L^{-1})\;\cong\;\bigoplus_{0\le r_1,\dots,r_\ell\le m-1} H^i\bigl(X,Q^{-1-r_1-\cdots-r_\ell}\bigr).
\]
In particular $H^i(X,Q^{-1})$ is a direct summand of $H^i(Y,L^{-1})$.
\end{theorem}

Thus, in every dimension and characteristic, a counterexample to Kodaira vanishing with a globally generated bundle
yields one with a very ample bundle. Raynaud's question for very ample bundles reduces to his question for globally
generated ones, and part (a) answers the latter on his own surface.

\paragraph{What is Raynaud's and what is added.} The curve, the ruled surface, the triple cover and the bundle $N$ are
those of \cite[\S\S1--3]{Raynaud}, with his Remarque~3 fixing the curve and $N$. The additions here are the following:
\begin{itemize}
\item the twist $Q=\OO_S(2T)\otimes f^*N$, with the double of the section $T$ in place of Raynaud's $E$, which makes the
bundle globally generated along $T$;
\item the computation of all its negative cohomology;
\item the very ampleness of $Q^{\otimes 9}$;
\item Theorem~\ref{thm:roots}.
\end{itemize}
Adjoining roots of every member of a basis of a linear system is a standard covering construction. Bhatt uses it to
make the pullback of an ample bundle globally generated \cite[proof of Proposition~7.2, p.~1771]{Bhatt}. The observation
that a failure of Kodaira vanishing survives a separable cover of degree prime to $p$, because $L^{-1}$ is a direct
summand of the pushforward, is made by Mukai \cite[p.~2]{Mukai}. What Theorem~\ref{thm:roots} adds is the very
ampleness. The construction is related to the Kummer covers of Bloch--Gieseker type, which make a pullback divisible.
Note that $Y$ is not the fibre product of $X$ with the Kummer map $[t_i]\mapsto[t_i^m]$ of projective space. That fibre
product is the quotient of $Y$ by the diagonal action of $\mu_m$, and the sections of $Q$ pulled back to $Y$ are what
separate the points of a diagonal orbit.

\paragraph{Context.} Ekedahl and Mukai showed that Kodaira vanishing for numerically positive (Mukai: nef and big)
line bundles holds on surfaces that are not of general type, apart essentially from Raynaud's examples
\cite[p.~99 and Theorem~II.1.3]{Ekedahl}\cite[Theorem~3]{Mukai}. Raynaud's surface in characteristic $2$ is
quasi-elliptic \cite[Remarque~1, pp.~277--278]{Raynaud}, so the surface $S$ of part (a) lies in that exception. For a
minimal surface of general type, Ekedahl proved $H^1(X,\omega_X^{-i})=0$ for all $i>0$, except possibly for $i=1$ in
one family in characteristic $2$ \cite[Main theorem~(i), p.~97]{Ekedahl}. So very ample counterexamples cannot be
found among pluricanonical bundles.
\begin{itemize}
\item Mukai \cite[Theorem~2]{Mukai} extended Raynaud's method to every dimension $n\ge2$, with ample bundles.
\item Zheng \cite{Zheng} refined and generalized Raynaud's surfaces of general type.
\item Bernasconi and Valloni recently showed that on infinitely many Raynaud surfaces with $p>2$ the $p$-th power of
the nonvanishing bundle is globally generated and separates points, restricting to a purely inseparable map of degree
$p$ on each fibre \cite[Theorem~1.2 and Propositions~3.13--3.14]{BV}. Their nonvanishing statement concerns the
bundle itself, not its $p$-th power, and very ampleness is not asserted.
\end{itemize}
None of these works contains a very ample surface counterexample. Nor did searches of Hexagon (23 November 2026)
find one.

\paragraph{Scope.} The example is in characteristic $2$. For odd $p$, Theorem~\ref{thm:roots} would give very ample
surface counterexamples as soon as a globally generated ample bundle with $H^1$ of its inverse nonzero is found on a
surface. Whether the twist used here has an analogue on Raynaud's odd-characteristic surfaces is not decided here.
The invariants of $Y$ are large, and no attempt is made to minimize them.

\section{Proof overview}

There are two difficulties. Raynaud's bundle $\mathcal I=\OO_X(E)\otimes N$ fails to be generated along the section
$E$. And even a generated bundle is far from very ample, which is what the question asks for.

\emph{Global generation.} The surface $S$ is a cyclic triple cover $\psi:S\to P$ of a ruled surface $\pi:P\to C$,
totally ramified along a section $E_\infty$ of $P$, with $\psi^*E_\infty=3T$. The pushforward
$\psi_*\OO_S(aT)$ splits into three eigensheaves, and the twist by $T$ shifts the pole order allowed along $E_\infty$
by $\lfloor (a+j)/3\rfloor$ in the $j$-th summand. With $a=1$, as for Raynaud's bundle, only the invariant summand has
sections, and they vanish on $T$. With $a=2$ the first eigensheaf acquires sections with a pole of exact order two along
$T$, which generate $\OO_S(2T)$ there. The same eigensheaf decomposition, together with relative duality on the
ruled surface, computes every $H^1(S,Q^{-n})$. Only $n=1$ contributes, through a trivial summand of $R^1$, so that
$h^1(Q^{-1})=1$.

\emph{Very ampleness.} Since $Q^{\otimes9}$ is pulled back from $P$, its very ampleness reduces to two standard
statements on the ruled surface: a line bundle is very ample, and its twist by one eigensheaf is generated. If $Q$ is
generated and $Q^{\otimes m}$ is very ample, let $t_i$ be an $m$-th root of a section $s_i$ of $Q^{\otimes m}$. On the
open set where a section $u$ of $Q$ does not vanish, the functions $s_i/u^m$ generate the coordinate ring, because the
$s_i$ form a basis of a very ample complete linear system. Upstairs the ratios $t_i/u$ are $m$-th roots of these
functions, so they generate the coordinate ring of the cover. Thus the roots $t_i$ together with the sections $u$ of
$Q$ embed $Y$. The cohomology of $L^{-1}$ is the sum, over all multi-indices of roots, of cohomology of negative powers
of $Q$ on $X$. Its summand of index zero is $H^1(Q^{-1})\neq0$, and all the others vanish for the surface $S$.

\emph{The hypotheses.} The characteristic $2$ is used in Raynaud's construction: the Artin--Schreier curve and the
inseparable bisection $C''$, whose equation $w^2+x=0$ is smooth because $d(w^2)=0$. The prime-to-$p$ degrees $3$ and $9$
make the triple cover and the root cover smooth and split their pushforwards into eigensheaves. Ampleness of $Q$ kills
the sections of all nontrivial summands, which makes $Y$ connected and fixes the numbers $562$ and $565$.

\section{Roots of a basis: proof of Theorem~\ref{thm:roots}}

This section proves the general construction; the surface is treated in the next two sections.

Write $Q^{-r}$ for $(Q^{\otimes r})^{-1}$. Since $m$ is invertible in $k$, adjoining an $m$-th root of a section $s$ of
$Q^{\otimes m}$ means taking $\Spec_X\bigoplus_{r=0}^{m-1}Q^{-r}$, with multiplication
$Q^{-r}\otimes Q^{-r'}\to Q^{-r-r'}$ when $r+r'<m$. When $r+r'\ge m$ the product is followed by multiplication by $s$,
from $Q^{-r-r'}$ to $Q^{-r-r'+m}$. Adjoining roots of $s_1,\dots,s_\ell$ simultaneously gives
\begin{equation}\label{eq:rootalgebra}
 \rho:Y=\Spec_X\Bigl(\bigotimes_{i=1}^\ell\bigoplus_{r=0}^{m-1}Q^{-r}\Bigr)\longrightarrow X,\qquad
 \rho_*\OO_Y=\bigoplus_{0\le r_i\le m-1}Q^{-(r_1+\cdots+r_\ell)} .
\end{equation}
This is finite and locally free of degree $m^\ell$, so $Y$ is projective.

\emph{Choice of the basis.} The tuples $(s_1,\dots,s_\ell)$ form an affine space over $k$. By Bertini's theorem for the
very ample bundle $Q^{\otimes m}$ over the algebraically closed field $k$ \cite[II.8.18]{Hartshorne}, applied
successively to $X$ and to the intersections already formed, the following conditions are nonempty and open:
\begin{itemize}
\item the $s_i$ form a basis;
\item each divisor $D_i=\{s_i=0\}$ is smooth;
\item the divisors meet transversally;
\item no $\dim X+1$ of them have a common point.
\end{itemize}
A finite intersection of nonempty open subsets of an irreducible variety is nonempty, and over the infinite field $k$
it has rational points. We take $(s_i)$ to be such a point.

\emph{Smoothness and connectedness.} At a point $x\in X$ the divisors $D_i$ through $x$ have local equations forming
part of a regular system of parameters. Adjoining $m$-th roots of these parameters gives again a regular local ring.
Adjoining $m$-th roots of the remaining $s_i$, which are units at $x$, is étale because $m$ is invertible. Hence $Y$ is
smooth. Since $Q^{\otimes m}$ is very ample, $Q$ is ample, so $H^0(X,Q^{-r})=0$ for $r>0$. By
\eqref{eq:rootalgebra}, $H^0(Y,\OO_Y)=H^0(X,\OO_X)=k$. A smooth proper scheme with $H^0(\OO)=k$ is connected, and
being smooth it is integral.

\emph{Cohomology.} By the projection formula, $\rho_*L^{-1}=Q^{-1}\otimes\rho_*\OO_Y$, and $\rho$ is finite. So
\eqref{eq:rootalgebra} gives the displayed decomposition of $H^i(Y,L^{-1})$, whose summand with all $r_i=0$ is
$H^i(X,Q^{-1})$.

\emph{Very ampleness.} On $Y$ let $t_i$ be the tautological section of $L$ with $t_i^m=\rho^*s_i$. It is the unit
section $1$ of the summand $\OO_X=Q\otimes Q^{-1}$ of $\rho_*L$, in the $i$-th factor with $r_i=1$. Let
$u_1,\dots,u_M$ be a basis of $H^0(X,Q)$. Consider the morphism $\phi:Y\to\PP^{\ell+M-1}$ given by the sections
$t_1,\dots,t_\ell,\rho^*u_1,\dots,\rho^*u_M$ of $L$; these have no common zero because the $u_j$ generate $Q$. Fix $j$,
put $u=u_j$, and let $X_u\subset X$ be the open set where $u\ne0$. Since $u^m$ is a section of the very ample bundle
$Q^{\otimes m}$ and the $s_i$ form a basis of its sections, $X_u$ is the standard affine chart of the embedding by
$|Q^{\otimes m}|$ given by $u^m\ne0$. Its coordinate ring is therefore generated over $k$ by the functions $s_i/u^m$. On
$\rho^{-1}(X_u)$ the bundle $L$ is trivialized by $u$. By \eqref{eq:rootalgebra},
\[
 \OO\bigl(\rho^{-1}(X_u)\bigr)=\OO(X_u)[z_1,\dots,z_\ell]/(z_i^m-s_i/u^m),\qquad z_i=t_i/u .
\]
Since $s_i/u^m=z_i^m$, this ring is generated over $k$ by the $z_i$ alone. Now $\rho^{-1}(X_u)$ is the inverse image under
$\phi$ of the standard chart $\{\rho^*u\ne0\}$ of the projective space, and the coordinate functions of that chart map
onto generators of its coordinate ring. So $\phi$ is a closed immersion over each such chart. The charts cover the image
because $Q$ is generated by the $u_j$. Thus $\phi$ is an immersion, and as $Y$ is proper it is a closed immersion. So
$L$ is very ample. \qed

\begin{remark}
The sections $\rho^*u_j$ cannot be omitted. The roots $t_i$ alone define the map to the fibre product of $X$ with the
Kummer map $[x_i]\mapsto[x_i^m]$ of $\PP^{\ell-1}$. That map identifies the points $(t_i)$ and $(\zeta t_i)$ for
$\zeta\in\mu_m$, so it is not injective.
\end{remark}

\section{Raynaud's surface in characteristic 2 and a generated bundle}

This section constructs $S$ and $Q$ and proves part (a) of Theorem~\ref{thm:main}. All objects are defined over $\F_2$,
and cohomology is computed after base change to $k$; dimensions do not change under field extension.

\subsection{The curve}
Let $C$ be the smooth projective model of $y^2+y=x^{15}$. The affine curve is smooth, since the partial derivative in
$y$ is $1$. The function $x$ is a coordinate at every affine point, so $dx$ has no affine zero, and differentiating gives
$dy=x^{14}dx$. A pole of $g^2+g$ has even order while $x^{15}$ has a pole of odd order at infinity. Hence $C$ is
geometrically integral, and the cover $C\to\PP^1$ has one point $\infty$ over $x=\infty$, with $v_\infty(x)=-2$ and
$v_\infty(y)=-15$.

Put $P_0=(0,0)$, $P_1=(0,1)$ and $t=y/x^8$. Then $v_\infty(t)=1$; at $P_0$ one has $v(y)=15$ and $v(x)=1$, so $v(t)=7$;
and at $P_1$, $v(t)=-8$. Dividing $y^2=y+x^{15}$ by $x^{16}$ gives $t^2=t x^{-8}+x^{-1}$. Multiplying by $x/t^2$, one
gets
\begin{equation}\label{eq:curve}
 \operatorname{div}(t)=7P_0+\infty-8P_1,\qquad x+t^{-2}=t^{-1}x^{-7},\qquad dt=x^{-8}dy=x^6dx .
\end{equation}
Since $t$ is a uniformizer at $\infty$, $v_\infty(dx)=v_\infty(dt)-6v_\infty(x)=12$. So $\operatorname{div}(dx)=12\infty$ and
$C$ has genus $7$. Its Weierstrass semigroup at $\infty$ contains the pole orders $2$ of $x$ and $15$ of $y$. Hence
$h^0(\OO_C(2\infty))=2$, with basis $1,x$, and $h^0(\OO_C(6\infty))=4$, with basis $1,x,x^2,x^3$. Indeed, a function with
a pole of order at most $6$ at $\infty$ and regular elsewhere is a polynomial in $x$ and $y$; since $y$ has pole order
$15$, it is a polynomial in $x$. Put $N=\OO_C(2\infty)$, a line bundle of degree $2$ generated by $1$ and $x$.

\subsection{The ruled surface and the branch curve}
On $U=C\setminus\{\infty\}$ and $V=C\setminus\{P_0,P_1\}$ define a rank-two bundle $E$ by frames $(e_0,e_1)$ and
$(v_0,v_1)$, related on $U\cap V$ by
\[
 v_0=e_0,\qquad v_1=t^{-6}(e_1-t^{-1}e_0).
\]
The section $e_0=v_0$ spans a trivial subbundle. The quotient is generated by $e_1$ on $U$ and by $t^{-6}e_1$ on $V$. Its
sections are the functions regular on $U$ with a pole of order at most $6$ at $\infty$, so the quotient is
$\OO_C(6\infty)=N^3$. This gives
\[
 0\to\OO_C\to E\to N^3\to0,\qquad\det E=N^3 .
\]
Consequently $\Sym^aE$ has a filtration whose successive quotients are $N^{3h}$, $0\le h\le a$.

Let $\pi:P=\PP_C(E)\to C$ be the bundle of rank-one quotients, so that $\pi_*\OO_P(j)=\Sym^jE$ for $j\ge0$. The global
section $e_0$ of $E$ defines a section $E_\infty=\{e_0=0\}$ of $\pi$, with $\OO_P(E_\infty)=\OO_P(1)$.

The quadratic expression $e_1^2+xe_0^2$ is a section of $\OO_P(2)$ over $U$. In characteristic $2$ one has
$e_1^2=t^{12}v_1^2+t^{-2}v_0^2$. By \eqref{eq:curve},
\[
 e_1^2+xe_0^2=t^{12}\bigl(v_1^2+c\,v_0^2\bigr),\qquad c=t^{-12}(x+t^{-2})=t^{-13}x^{-7}.
\]
The function $c$ is regular on $V$ and has valuation $-13+14=1$ at $\infty$. So the two local expressions define a curve
$C''\subset P$, with $\OO_P(C'')=\OO_P(2)\otimes\pi^*N^{-6}$ (the factor $t^{12}$ matches $N^{-6}=\OO_C(-12\infty)$).
On $E_\infty$ one has $e_0=0$, where the expression is $e_1^2\ne0$; so $C''$ and $E_\infty$ are disjoint. Away from
$\infty$, in the coordinate $w=e_1/e_0$, the curve $C''$ is $w^2+x=0$, whose differential $dx$ does not vanish. Near
$\infty$ its equation $v^2+c$ has differential $dc\ne0$, since $c$ is a uniformizer there. Thus $C''$ is smooth.

\subsection{The triple cover}
Set $M=\OO_P(-1)\otimes\pi^*N^2$. Then $M^{-3}=\OO_P(3)\otimes\pi^*N^{-6}=\OO_P(E_\infty+C'')$. The section $e_0\cdot(e_1^2+xe_0^2)$ of
$M^{-3}$ vanishes on $E_\infty+C''$ and defines a cyclic triple cover
\[
 \psi:S=\Spec_P(\OO_P\oplus M\oplus M^2)\longrightarrow P,
\]
with multiplication $M\otimes M^2\to\OO_P$ given by that section. The branch divisor is smooth, and $3$ is invertible.
Locally the cover is $z^3=u$ along a branch equation $u$, and étale elsewhere. So $S$ is smooth and projective. On the
chart $e_0\ne0$ it is
\[
 y^2+y=x^{15},\qquad z^3=w^2+x .
\]
This is the surface of Raynaud's \S2 with the data of his Remarque~3 \cite[pp.~275--278]{Raynaud}. Since $M$ and $M^2$ have negative degree on the
fibres of $\pi$, $H^0(S,\OO_S)=k$; so $S$ is connected, hence integral. Write $f=\pi\psi$, and let $T$ be the reduced
inverse image of $E_\infty$. Then $\psi^*E_\infty=3T$ and $T\simeq C$.

\begin{lemma}\label{lem:floor}
For every integer $a$,
$\displaystyle\psi_*\OO_S(aT)=\bigoplus_{j=0}^2M^j\bigl(\lfloor (a+j)/3\rfloor E_\infty\bigr)$.
\end{lemma}
\begin{proof}
Near $E_\infty$ let $u$ be a local equation of $E_\infty$ and $z$ a local generator of $M$ in the algebra, so that
$z^3=u\cdot(\text{unit})$, $v_T(z)=1$ and $v_T(u)=3$. A local section of $\psi_*\OO_S$ with poles along $E_\infty$ is
$\sum_j\varphi_jz^j$, with $\varphi_j$ rational functions on $P$. The valuations $v_T(\varphi_jz^j)=3v_{E_\infty}(\varphi_j)+j$
are distinct modulo $3$, so the sum has valuation at least $-a$ exactly when each term has. That is,
$-v_{E_\infty}(\varphi_j)\le(a+j)/3$.
\end{proof}

\begin{proposition}\label{prop:Q}
The line bundle $Q=\OO_S(2T)\otimes f^*N$ is ample and globally generated, and
\[
 h^0(S,Q)=6,\quad h^1(S,Q^{-1})=1,\quad H^1(S,Q^{-n})=0\ (n\ge2),\quad h^1(S,\OO_S)=9 .
\]
\end{proposition}
\begin{proof}
\emph{Sections and generation.} By Lemma~\ref{lem:floor} with $a=2$, and $\OO_P(E_\infty)=\OO_P(1)$, the three summands of
$\psi_*Q$ are $N$, $M(E_\infty)\otimes N=\pi^*N^3$, and $M^2(E_\infty)\otimes N=\OO_P(-1)\otimes\pi^*N^5$. The last has no
direct images. Hence $f_*Q=N\oplus N^3$, and $h^0(S,Q)=2+4=6$. On the chart above the sections are
$1,x,z,zx,zx^2,zx^3$.

The sections from the summand $N$ are regular functions, generated by $1$ and $x$. They generate $Q$ away from $T$,
where $\OO_S(2T)$ is trivialized by $1$. The sections from the second summand have the form $\varphi z$, with $\varphi$
having a simple pole along $E_\infty$, so they have valuation $1-3=-2$ along $T$. Since $N^3$ is generated by its
sections, they generate $\OO_S(2T)\otimes f^*N$ along $T$. So $Q$ is globally generated. It is ample, because its ninth
power is very ample (Proposition~\ref{prop:Q9}).

\emph{Negative powers.} For $n\ge1$, Lemma~\ref{lem:floor} gives the $j$-th summand of $\psi_*Q^{-n}$ as
\[
 \OO_P(a_j)\otimes\pi^*N^{2j-n},\qquad a_j=\lfloor(-2n+j)/3\rfloor-j\le-1 .
\]
So $\pi_*$ of every summand vanishes, and the Leray spectral sequence gives $H^1(S,Q^{-n})=H^0(C,R^1\pi_*\psi_*Q^{-n})$.
Relative duality for $\pi$, with $\omega_{P/C}=\OO_P(-2)\otimes\pi^*\det E$, gives
\[
 R^1\pi_*\OO_P(a)\cong(\Sym^{-a-2}E)^\vee\otimes N^{-3}\quad(a\le-2),\qquad R^1\pi_*\OO_P(-1)=0 .
\]
For $n=1$ one has $(a_0,a_1,a_2)=(-1,-2,-2)$, and the nonzero direct images are $N^{-3}\otimes N=N^{-2}$ (for $j=1$) and
$N^{-3}\otimes N^3=\OO_C$ (for $j=2$). Hence $h^1(S,Q^{-1})=h^0(N^{-2})+h^0(\OO_C)=0+1=1$.

For $n\ge2$, dualizing the filtration of $\Sym^{-a_j-2}E$ shows that every nonzero summand of $R^1\pi_*\psi_*Q^{-n}$ is
a successive extension of the line bundles $N^{2j-n-3-3h}$, $h\ge0$. Their exponents are at most $2\cdot2-n-3=1-n<0$,
so these bundles have negative degree. A successive extension of bundles without sections has none, so
$H^1(S,Q^{-n})=0$.

\emph{Irregularity.} The summands of $\psi_*\OO_S$ are $\OO_P$, $\OO_P(-1)\otimes\pi^*N^2$ and $\OO_P(-2)\otimes\pi^*N^4$,
with $\pi_*$ equal to $\OO_C,0,0$ and $R^1\pi_*$ equal to $0,0,N$. As $C$ is a curve, Leray gives
$h^1(S,\OO_S)=h^1(\OO_C)+h^0(N)=7+2=9$.
\end{proof}

\section{The ninth power is very ample}

This section proves that $Q^{\otimes9}$ is very ample and counts its sections, the input of Theorem~\ref{thm:roots}.
Put $B=\OO_P(6)\otimes\pi^*N^9$. Since $18T=\psi^*(6E_\infty)$, one has $Q^{\otimes9}=\psi^*B$. Its pushforward is
$B\oplus BM\oplus BM^2$, whose three summands are $\OO_P(6-i)\otimes\pi^*N^{9+2i}$ for $0\le i\le2$.

\begin{proposition}\label{prop:Q9}
The bundle $Q^{\otimes9}$ is very ample on $S$ and $h^0(S,Q^{\otimes9})=210+186+160=556$.
\end{proposition}
\begin{proof}
For $(a,b)\in\{(6,9),(5,11),(4,13)\}$, the bundle $\Sym^aE\otimes N^b=\pi_*(\OO_P(a)\otimes\pi^*N^b)$ is filtered with
quotients $N^{b+3h}$, $0\le h\le a$, of degree $2b+6h\ge18$. After twisting down by any effective divisor of degree
two on $C$ the degrees are still at least $16>2g-2=12$. By Serre duality these line bundles, and hence (by the
filtration) the vector bundles, have vanishing $H^1$ after such a twist. Their sections therefore surject onto every
fibre over a subscheme of $C$ of length at most two. By Riemann--Roch each quotient of degree $d$ contributes $d-6$
sections, so
\[
 h^0(B)=\sum_{h=0}^6(12+6h)=210,\qquad h^0(BM)=\sum_{h=0}^5(16+6h)=186,
\]
\[
 h^0(BM^2)=\sum_{h=0}^4(20+6h)=160 .
\]

$B$ is very ample on $P$. Let $Z\subset P$ be a subscheme of length two and $W=\pi(Z)$. The restriction
$H^0(P,B)=H^0(C,\Sym^6E\otimes N^9)\to H^0(\pi^{-1}W,B)$ is surjective by the preceding paragraph. On $\pi^{-1}W$,
which is a $\PP^1$-bundle over $W$, the bundle $B$ is $\OO(6)$. Its sections separate the length-two subscheme $Z$:
\begin{itemize}
\item if $W$ is one point, $Z$ lies in one fibre $\PP^1$, and $\OO(6)$ is very ample there;
\item otherwise $Z\to W$ is an isomorphism, and $Z$ is the image of a section over $W$.
\end{itemize}
The same surjectivity onto single fibres shows that $BM=\OO_P(5)\otimes\pi^*N^{11}$ is globally generated.

Now let $b\in H^0(P,B)$ and $D(b)=\{b\ne0\}$, a standard affine chart of the embedding by $B$. The ring $\OO(D(b))$ is
generated over $k$ by the ratios $s/b$, $s\in H^0(P,B)$. Over $D(b)$ the algebra $\psi_*\OO_S$ is
$\OO\oplus M\oplus M^2$, generated as an $\OO(D(b))$-algebra by $H^0(D(b),M)$, because $M\otimes M\to M^2$ is
surjective. Since $BM$ is globally generated, the sheaf $M|_{D(b)}=(BM\otimes B^{-1})|_{D(b)}$ is generated by the
ratios $\sigma/b$ with $\sigma\in H^0(P,BM)$. On the affine $D(b)$ these ratios therefore generate the module
$H^0(D(b),M)$. All these ratios are ratios of sections of $\psi^*B$, so the map defined by the complete linear system
of $\psi^*B$ is a closed immersion over each chart $\{b\ne0\}$. These charts cover the image, since $B$ is generated.
As $S$ is proper, $\psi^*B=Q^{\otimes 9}$ is very ample.
\end{proof}

\section{Proof of Theorem~\ref{thm:main}}

This section assembles the parts. Part (a) is Proposition~\ref{prop:Q}, with ampleness from
Proposition~\ref{prop:Q9}.

For part (b), apply Theorem~\ref{thm:roots} to $X=S$ and $Q$ with $m=9$, which is prime to $2$. Here $Q$ is generated
by Proposition~\ref{prop:Q} and $Q^{\otimes9}$ is very ample with $\ell=556$ sections by Proposition~\ref{prop:Q9}. We
obtain a smooth projective connected surface $Y$, finite and flat of degree $9^{556}$ over $S$, with $L=\rho^*Q$ very
ample. By \eqref{eq:rootalgebra} the invariants are sums over the multi-indices $r=(r_i)$, with $|r|=\sum r_i$:
\begin{itemize}
\item $H^0(Y,L)=\bigoplus_rH^0(S,Q^{1-|r|})$. The index $|r|=0$ contributes $h^0(Q)=6$. Each of the $556$ indices with
$|r|=1$ contributes $h^0(\OO_S)=1$. Larger $|r|$ give negative powers of the ample $Q$, which have no sections. So
$h^0(L)=562$.
\item $H^1(Y,L^{-1})=\bigoplus_rH^1(S,Q^{-1-|r|})$. Only $|r|=0$ contributes, by Proposition~\ref{prop:Q}. So
$h^1(L^{-1})=1$.
\item $H^1(Y,\OO_Y)=\bigoplus_rH^1(S,Q^{-|r|})$. This is $9$ from $|r|=0$ plus $556\cdot1$ from $|r|=1$, the others
vanishing. So $h^1(\OO_Y)=565$.
\end{itemize}
The sections $t_i$ (the $556$ indices with $|r|=1$) and the six pullbacks of sections of $Q$ form a basis of $H^0(Y,L)$.
So the embedding of Theorem~\ref{thm:roots} is by the complete linear system $|L|$.

Finally let $H\in|L|$ be smooth and connected. The sequence $0\to L^{-1}\to\OO_Y\to\OO_H\to0$ and the isomorphism
$H^0(\OO_Y)=k=H^0(\OO_H)$ identify $H^1(Y,L^{-1})$ with the kernel of $H^1(Y,\OO_Y)\to H^1(H,\OO_H)$. This is the map
induced by restriction on the tangent spaces at the identity of the Picard schemes. \qed

\section{What remains open}

\begin{itemize}
\item \emph{Odd characteristic.} By Theorem~\ref{thm:roots}, a very ample counterexample on a surface in characteristic
$p$ exists as soon as a globally generated ample $Q$ with $H^1(Q^{-1})\ne0$ exists on some surface. For $p$ odd we have
not checked whether a twist of Raynaud's bundle along the cusp section is globally generated.
\item \emph{Size.} The surface $Y$ has irregularity $565$ and its embedding lies in $\PP^{561}$. Minimizing the degree
of the cover, the irregularity or the embedding dimension is not attempted.
\item \emph{Explicit equations.} The proof shows that $Y$ can be defined over a finite field, but gives neither a
specific field nor a specific basis $(s_i)$.
\item \emph{Picard kernels.} Whether some smooth surface has a very ample (or merely ample) $L$ and a smooth
$H\in|L|$ for which $\Pic(Y)\to\Pic(H)$ has a nontrivial kernel on geometric points is not touched by this
construction. It controls tangent spaces only.
\end{itemize}

\section*{Method and checks}

The construction and its proofs were produced by GPT-6 Astra (ChatGPT Pro mode). That run checked them by a
reconstruction from the statement alone by a separate instance of the same model, and by two computer recomputations
of the dimensions $210$, $186$ and $160$. Claude Opus 5.5 at effort high made a targeted check of every step, recomputing
every invariant by hand. It also wrote this text: it separated Theorem~\ref{thm:roots} from the example, wrote out the
chart argument for Proposition~\ref{prop:Q9}, and replaced the run's ground field (an algebraic closure of a purely
transcendental extension of $\F_2$) by an arbitrary algebraically closed field of characteristic $2$, with the Bertini
argument of Section~3. No human has checked the mathematics.

The editor read at the source Raynaud's article \cite{Raynaud} (all of pp.~273--278), the opening pages of
\cite{LauritzenRao}, \cite{Mukai} and \cite{Ekedahl}, and the statements of \cite{BV} cited above. The locators in \cite{Bhatt} and \cite{Zheng} are those recorded by the run, and \cite{Lauritzen} is cited for its
result as \cite{LauritzenRao} describe it. The journal data of \cite{Lauritzen}, \cite{LauritzenRao} and \cite{Mukai} were
checked against bibliographic records on 23 November 2026. Our text is dedicated to the public domain under CC0~1.0. The records of the checks and the programs are at the
permanent page of this work, WB-0003: \url{\workbenchurl}.

\begin{thebibliography}{99}\small
\bibitem{BV} F.~Bernasconi and D.~Valloni, \emph{Brauer--Manin obstruction from failure of Kodaira vanishing in positive
characteristic}, arXiv:2605.28278 (2026).
\bibitem{Bhatt} B.~Bhatt, \emph{Derived splinters in positive characteristic}, Compos. Math. 148 (2012), 1757--1786.
\bibitem{Ekedahl} T.~Ekedahl, \emph{Canonical models of surfaces of general type in positive characteristic}, Publ. Math.
IHES 67 (1988), 97--144.
\bibitem{Hartshorne} R.~Hartshorne, \emph{Algebraic Geometry}, Springer, 1977.
\bibitem{Lauritzen} N.~Lauritzen, \emph{Embeddings of homogeneous spaces in prime characteristics}, Amer. J. Math. 118
(1996), 377--387; arXiv:alg-geom/9502016.
\bibitem{LauritzenRao} N.~Lauritzen and A.~P.~Rao, \emph{Elementary counterexamples to Kodaira vanishing in prime
characteristic}, Proc. Indian Acad. Sci. Math. Sci. 107 (1997), 21--25; arXiv:alg-geom/9604012.
\bibitem{Mukai} S.~Mukai, \emph{Counterexamples to Kodaira's vanishing and Yau's inequality in positive
characteristics}, Kyoto J. Math. 53 (2013), 515--532; preprint RIMS-1736 (2011).
\bibitem{Raynaud} M.~Raynaud, \emph{Contre-exemple au ``vanishing theorem'' en caract\'eristique $p>0$}, in
C.~P.~Ramanujam---a tribute, Tata Inst. Fund. Res. Studies in Math. 8, Springer, 1978, 273--278.
\bibitem{Zheng} X.~Zheng, \emph{Counterexamples of Kodaira vanishing for smooth surfaces of general type in positive
characteristic}, J. Pure Appl. Algebra 221 (2017), 2431--2444.
\end{thebibliography}
\end{document}
